% Irodov Part III §3.1 original 28--40, PDF pp.104--106,
% printed pp.95--97, answer pages 306--307.
\electroproblem{28}{95--96}
一个均匀体电荷密度为 $\rho$ 的实心球内，
挖去一个球形空腔。
设空腔中心相对原球心的位移矢量为 $\bm a$。
若介电常数为1，求空腔内部电场。
\exsolution
用叠加原理：
“挖空”相当于
完整均匀带电球（密度 $\rho$）
与一颗在空腔位置、密度为 $-\rho$ 的球叠加。
在均匀带电球内部电场为
$\bm E=\rho\bm r/(3\varepsilon_0)$。
空腔内任一点的原球心位置为 $\bm r$，
相对空腔中心位置为 $\bm r-\bm a$，
故
\[
\bm E_{\rm cavity}
=\frac{\rho\bm r}{3\varepsilon_0}
-\frac{\rho(\bm r-\bm a)}{3\varepsilon_0}
=\boxed{\frac{\rho\bm a}{3\varepsilon_0}}.
\]
\excheck 空腔场为\emph{均匀场}，
与腔半径及腔内观察点位置无关；
$\bm a$ 按原球心指向腔心定义。

\electroproblem{29}{96}
均匀带电的无限长实心圆柱，
体电荷密度为 $\rho$。
内部挖有一条平行轴线的圆柱形空腔，
两轴相隔距离 $a$。
介电常数为1，求空腔内电场。
\exsolution
以截面位置矢量表示，
均匀带电无限长圆柱\emph{内部}电场为
$\bm E=\rho\bm r_\perp/(2\varepsilon_0)$。
挖空用密度为 $-\rho$ 的小圆柱叠加。
设位移 $\bm a$ 从原圆柱轴指向空腔轴，
则空腔内一点对小圆柱轴坐标为
$\bm r_\perp-\bm a$，
所以
\[
\boxed{\bm E_{\rm cavity}=
\frac{\rho\bm r_\perp}{2\varepsilon_0}
-\frac{\rho(\bm r_\perp-\bm a)}{2\varepsilon_0}
=\frac{\rho\bm a}{2\varepsilon_0}}.
\]
\excheck 空腔场与位置无关；
因圆柱体是二维横截面高斯对称，
系数为 $1/2$，
而球形空腔第28题系数为 $1/3$。

\electroproblem{30}{96}
两个同轴细圆环，半径相同为 $R$，
相距 $a$，分别带电 $+q,-q$。
求两个圆环中心之间的电势差。
\exsolution
细圆环总电荷为 $Q$ 时，
距环心轴向距离 $z$ 处电势为
$\varphi(z)=kQ/\sqrt{R^2+z^2}$。
设正环中心为 A、负环中心为 B，
则
\[
\varphi_A=k\left(\frac qR-\frac q{\sqrt{R^2+a^2}}\right),
\quad
\varphi_B=k\left(-\frac qR+\frac q{\sqrt{R^2+a^2}}\right).
\]
因此
\[
\boxed{\varphi_A-\varphi_B=
\frac q{2\pi\varepsilon_0R}
\left(1-\frac1{\sqrt{1+(a/R)^2}}\right)}.
\]
\excheck 当 $a\to0$，两个正负圆环重合，
电势差趋零；$a\to\infty$ 时跨中心电势差
趋 $q/(2\pi\varepsilon_0R)$。

\electroproblem{31}{96}
一根无限长均匀线电荷，线密度
$\lambda=0.40\,\mathrm{\mu C/m}$。
点2到导线的距离是点1的 $\eta=2.0$ 倍。
求两点电势差 $\varphi_1-\varphi_2$。
\exsolution
由高斯定律，径向电场
$E_r=\lambda/(2\pi\varepsilon_0r)$。
电势差为电场线积分的相反数：
\[
\varphi_1-\varphi_2=
\int_{r_1}^{r_2}E_r\,\mathrm dr
=\frac{\lambda}{2\pi\varepsilon_0}
\ln\frac{r_2}{r_1}
=\frac{\lambda}{2\pi\varepsilon_0}\ln\eta.
\]
代入题值：
\[
\boxed{\varphi_1-\varphi_2\simeq
5.0\times10^3\,\mathrm V=5.0\,\mathrm{kV}}.
\]
\excheck 无限线电荷\emph{绝对电势}
不能用无穷远为零来归一化，
但有限两点电势差有意义。

\electroproblem{32}{96}
半径 $R$ 的半球壳均匀面带电密度
$\sigma$。求球心处电势和电场模。
\exsolution
壳面上各点距球心都为 $R$，
半球面积 $2\pi R^2$，
故电势
\[
\varphi(0)=k\int_{\rm hemi}\frac{\sigma\,\mathrm dS}R
=k\frac{\sigma(2\pi R^2)}R
=\boxed{\frac{\sigma R}{2\varepsilon_0}}.
\]
设半球位于 $z>0$，
中心处场
$\mathrm d\bm E=-k\sigma\bm n\,\mathrm d\Omega$。
其轴向积分
$\int_{\rm hemi}\cos\theta\,\mathrm d\Omega
=2\pi\int_0^{\pi/2}\cos\theta\sin\theta\,\mathrm d\theta=\pi$，
因此
\[
\boxed{\bm E(0)=-\frac{\sigma}{4\varepsilon_0}\bm e_z},
\quad
\boxed{|\bm E|=\frac{|\sigma|}{4\varepsilon_0}}.
\]
\excheck 方向指向与带正电的半球相反一侧，
电势始终标量相加，无方向抵消。

\electroproblem{33}{96}
半径 $R$ 的无限薄均匀带电圆盘，
面电荷密度为 $\sigma$。
求沿圆盘轴线、距中心 $z>0$ 处的电势和场强；
分析 $z\to0^+$ 与 $z\gg R$。
\exsolution
把圆盘分解为半径 $s$、宽度 $\mathrm ds$
的同心圆环，
环电荷
$\mathrm dq=2\pi\sigma s\,\mathrm ds$。
轴上电势积分：
\[
\varphi(z)=k\int_0^R
\frac{2\pi\sigma s\,\mathrm ds}{\sqrt{z^2+s^2}}
=\boxed{\frac{\sigma}{2\varepsilon_0}
(\sqrt{z^2+R^2}-z)}.
\]
由 $\bm E=-\nabla\varphi$，
\[
\boxed{E_z(z)=\frac{\sigma}{2\varepsilon_0}
\left(1-\frac{z}{\sqrt{z^2+R^2}}\right)}.
\]
当 $z\to0^+$：
$\varphi\to\sigma R/(2\varepsilon_0)$、
$E_z\to\sigma/(2\varepsilon_0)$。
当 $z\gg R$，
$\sqrt{z^2+R^2}-z\simeq R^2/(2z)$；
设 $q=\pi R^2\sigma$，
\[
\boxed{\varphi\simeq\frac{q}{4\pi\varepsilon_0z},
\quad E_z\simeq\frac{q}{4\pi\varepsilon_0z^2}}.
\]
\excheck 在圆盘另一侧 $z<0$，
电势表达式须用 $|z|$，
电场法向方向随 $z$ 改变。

\electroproblem{34}{96}
半径为 $R$ 的薄圆盘均匀面带电密度
$\sigma$。求圆盘\emph{边缘}处电势
（无穷远电势取零）。
\exsolution
以圆盘边缘上的观察点作为极坐标原点。
取极轴指向圆盘中心，
一条与该方向成角 $\theta$ 的射线，
在圆盘内部的最长半径是
$s_{\max}=2R\cos\theta$，
$-\pi/2\le\theta\le\pi/2$。
从观察点到面电荷元的距离是 $s$，
面积元
$\mathrm dA=s\,\mathrm ds\,\mathrm d\theta$。
因此
\[
\varphi=k\sigma
\int_{-\pi/2}^{\pi/2}
\int_0^{2R\cos\theta}
\frac{s\,\mathrm ds\,\mathrm d\theta}{s}
=k\sigma(2R)\int_{-\pi/2}^{\pi/2}
\cos\theta\,\mathrm d\theta.
\]
积分后
\[
\boxed{\varphi_{\rm rim}=
4k\sigma R=\frac{\sigma R}{\pi\varepsilon_0}}.
\]
\excheck 圆盘边缘与中心不同；
中心电势是 $\sigma R/(2\varepsilon_0)$，
边缘电势较小，两者比值为 $2/\pi$。

\electroproblem{35}{96}
静电势为 $\varphi(\bm r)=\bm a\cdot\bm r$，
其中 $\bm a$ 为常矢量。
求电场强度。
\exsolution
写出笛卡尔分量
$\varphi=a_xx+a_yy+a_zz$，
其梯度
$\nabla\varphi=(a_x,a_y,a_z)=\bm a$。
由电势—电场关系
$\bm E=-\nabla\varphi$，
直接得到
\[
\boxed{\bm E=-\bm a}.
\]
\excheck 电场处处均匀，
方向与电势增长最快的方向相反。

\electroproblem{36}{96}
在 $(x,y)$ 平面中给定静电势：
(a) $\varphi=a(x^2-y^2)$；
(b) $\varphi=axy$。
求 $\bm E(x,y)$，并说明电场线的大致形状。
\exsolution
(a) 对电势求梯度：
\[
\boxed{\bm E_a=
-2ax\,\bm i+2ay\,\bm j}.
\]
场线微分方程
$\mathrm dy/\mathrm dx=E_y/E_x=-y/x$，
积分为
\[
\boxed{xy=\text{常数}}
\]
（坐标轴是分界线）。
(b) 同理
\[
\boxed{\bm E_b=-ay\,\bm i-ax\,\bm j}.
\]
场线
$\mathrm dy/\mathrm dx=x/y$，
积分为
\[
\boxed{y^2-x^2=\text{常数}}.
\]
两组场线都是双曲线族，
分别以坐标轴和斜线 $y=\pm x$ 为渐近线，
箭头方向由 $\bm E$ 的具体分量决定。
\excheck \textbf{本版书末第36(b)题符号疑误：}
简答将 $E_y$ 印作 $+ax$，
但由 $E_y=-\partial_y(axy)=-ax$ 必为负号。
若用书末符号，则
$\partial_xE_y-\partial_yE_x\ne0$，
也不符合静电场无旋条件。

\electroproblem{37}{96}
电势
$\varphi=a(x^2+y^2)+bz^2$。
求电场矢量及模，并讨论等势面
(a) $a>0,b>0$；
(b) $a>0,b<0$ 的形状。
\exsolution
直接对三坐标求负梯度：
\[
\boxed{\bm E=
-2a(x\bm i+y\bm j)-2bz\bm k},
\]
\[
\boxed{|\bm E|=
2\sqrt{a^2(x^2+y^2)+b^2z^2}}.
\]
等势面满足
$a(x^2+y^2)+bz^2=\varphi_*$。
(a) $a,b>0$ 且 $\varphi_*>0$，
是绕 $z$ 轴的旋转椭球面，
半轴 $\sqrt{\varphi_*/a}$、
$\sqrt{\varphi_*/b}$。
(b) $a>0,b<0$，
当 $\varphi_*>0$ 是单叶双曲面，
$\varphi_*=0$ 是圆锥面，
$\varphi_*<0$ 是双叶双曲面。
\excheck 等势面的分类依赖等势常数 $\varphi_*$
的符号，不可只依赖 $a,b$。

\electroproblem{38}{96}
半径 $R$ 的实心球内均匀分布总电荷 $q$，
球内外介电常数均为1。
求 (a) 球心电势；
(b) 距球心 $r<R$ 处电势。
\exsolution
由高斯定律，
球内场
$E(r)=qr/(4\pi\varepsilon_0R^3)$，
球外场
$E(r)=q/(4\pi\varepsilon_0r^2)$。
以无穷远电势为零，
球面电势
$\varphi(R)=q/(4\pi\varepsilon_0R)$。
球内
\[
\varphi(r)=\varphi(R)+\int_r^R E(s)\,\mathrm ds
=\frac{q}{4\pi\varepsilon_0R}
+\frac{q}{8\pi\varepsilon_0R^3}(R^2-r^2).
\]
因此
\[
\boxed{\varphi(r)=
\frac{q}{8\pi\varepsilon_0R}
\left(3-\frac{r^2}{R^2}\right)}.
\]
在球心 $r=0$：
\[
\boxed{\varphi(0)=\frac{3q}{8\pi\varepsilon_0R}}.
\]
\excheck 在球面 $r=R$ 与外部库仑电势连续，
且中心场 $-\varphi'(0)=0$。

\electroproblem{39}{97}
证明点偶极矩 $\bm p$ 在远场的电势为
$\varphi=k(\bm p\cdot\bm r)/r^3$，
并求电场模作为 $r,\theta$ 的函数，
其中 $\theta$ 是 $\bm p$ 与 $\bm r$ 夹角。
\exsolution
把偶极子视为两等量异号电荷
$+q,-q$，间距矢量为 $\bm d$，
令 $q|\bm d|\to p$ 而 $|\bm d|\to0$。
电势
\[
\varphi=kq\left(
\frac1{|\bm r-\bm d/2|}
-\frac1{|\bm r+\bm d/2|}\right)
\simeq kq\frac{\bm d\cdot\bm r}{r^3}
=\boxed{\frac{\bm p\cdot\bm r}
{4\pi\varepsilon_0r^3}}.
\]
球坐标中
$\varphi=kp\cos\theta/r^2$。
因此
$E_r=-\partial_r\varphi
=2kp\cos\theta/r^3$，
$E_\theta=-(1/r)\partial_\theta\varphi
=kp\sin\theta/r^3$。
合场模
\[
\boxed{E=\frac{|p|}{4\pi\varepsilon_0r^3}
\sqrt{1+3\cos^2\theta}}.
\]
\excheck 轴线上 $\theta=0$ 得
$E=2kp/r^3$；
赤道面 $\theta=\pi/2$ 得 $E=kp/r^3$，
轴上场为赤道场的两倍。

\electroproblem{40}{97}
点偶极矩 $\bm p=p\bm e_z$ 位于原点，
求电场沿 $z$ 轴的分量 $E_z$
及垂直 $z$ 轴方向的分量 $E_\perp$；
在何种角度电场垂直于 $\bm p$？
\exsolution
偶极场矢量形式
\[
\bm E=k\frac{3(\bm p\cdot\bm n)\bm n-\bm p}{r^3},
\qquad\bm n=\bm r/r.
\]
用
$\bm n=\sin\theta\,\bm e_\perp+
\cos\theta\,\bm e_z$ 分解：
\[
\boxed{E_z=\frac{p}{4\pi\varepsilon_0r^3}
(3\cos^2\theta-1)},
\]
\[
\boxed{E_\perp=\frac{3p}{4\pi\varepsilon_0r^3}
\sin\theta\cos\theta}.
\]
电场与 $\bm p$ 垂直等价于 $E_z=0$，
所以
\[
\cos^2\theta=\frac13,\qquad
\boxed{\theta=\arccos(1/\sqrt3)\simeq54.7^\circ
\quad\text{或}\quad125.3^\circ}.
\]
\excheck 本版书末第二角印约 $123.5^\circ$，
但应严格是
$180^\circ-54.7356^\circ=125.2644^\circ$；
可能存在数字排印错误。
